\section{A nonelliptic compact family on the same lattice}
\label{sec:v37-support-family}
The action principle is not confined to quadrics or to centrally
symmetric obstacles. The following family provides a separate
verification on the same triangular lattice, with no change in the
Euclidean reflection law or in the observation cell.

Let $n_\psi=(\cos\psi,\sin\psi)$, $t_\psi=(-\sin\psi,\cos\psi)$.
For parameters $\lambda=(R,e,f,\theta,\phi)$ set
\begin{align}\label{eq:v37-support-family}
 h_\lambda(\psi)&=R+e\cos(3(\psi-\theta))
                         +f\cos(4(\psi-\phi)),\notag\\
 R&\in[91/200,93/200],\qquad |e|+|f|\le1/250,\qquad
                       \theta,\phi\in\R/2\pi\Z.
\end{align}
Define the obstacle by its support function $h_\lambda$, equivalently
by the boundary parametrization
$q_\lambda(\psi)=h_\lambda(\psi)n_\psi+h_\lambda'(\psi)t_\psi$.
Place its identical translates at every lattice point.

\begin{proposition}[Verification of the nonelliptic family]
\label{prop:v37-support-family}
The family \eqref{eq:v37-support-family} satisfies the separated
hypotheses of Section~\ref{sec:compact-family-local-principle}.
Its stationary microscopic local law, regular endpoint-selected law,
and normalized positive-likelihood comparison hold uniformly in
$\lambda$. The local laws are qualitative $o(1)$ statements.
The obstacle perimeter and area are
\begin{equation}\label{eq:v37-support-area}
 L_\lambda=2\pi R,\qquad
 |O_\lambda|=\pi R^2-4\pi e^2-\tfrac{15}{2}\pi f^2,
 \qquad
 \bar\tau_\lambda=\frac{\pi(\sqrt3/2-|O_\lambda|)}{L_\lambda}.
\end{equation}
The physical covariance is continuous and uniformly positive definite.
If $e\ne0$, the obstacle is not centrally symmetric and hence is not
an ellipse.
\end{proposition}
\begin{proof}
The support function and curvature radius satisfy
\begin{equation}\label{eq:v37-support-bounds}
 \frac{451}{1000}\le h_\lambda\le\frac{469}{1000},\qquad
 \frac{79}{200}\le h_\lambda+h_\lambda''
                                  \le\frac{21}{40}.
\end{equation}
The second bound follows from
$h+h''=R-8e\cos(3(\psi-\theta))-15f\cos(4(\psi-\phi))$.
In particular $q_\lambda'=(h+h'')t_\psi$ never vanishes. Positive
curvature makes this a regular strictly convex closed boundary with
the stated support function. Its curvature lies between $40/21$ and
$200/79$. All boundary derivatives of any fixed order are uniformly
bounded. The normalized arclength change is smooth with uniformly
bounded derivatives and inverse derivatives of every fixed order.

The support inequalities imply
$B(0,451/1000)\subset O_\lambda\subset B(0,469/1000)$.
Different obstacles thus have a gap at least $31/500$.
Containment of the radius-$0.45$ disks transfers their uniform finite
horizon to the present tables. The transverse finite-horizon condition
and the finite-cover mixing required in the arithmetic follow from
Proposition~\ref{prop:v37-finite-cover-mixing}. The unweighted common
spaces and normalized Jacobian estimates follow locally from the
uniform smooth dispersing geometry and its compact parameter cover.
The one-flight action identity uses arclength and is unchanged by the
support parametrization.

Here are the additional partition checks; smooth convexity alone is
not substituted for them. There are finitely many possible neighboring
obstacles within the uniform horizon. On bounded rational trigonometric
charts the coordinates of $q_\lambda$, the velocity, and the
intersection equation with each neighboring translated boundary are
semialgebraic. The coefficients $e,f,\cos(j\theta),\sin(j\theta)$
and their $\phi$ counterparts range in a fixed compact algebraic
parameter set. Selecting the least admissible positive intersection
uses only finitely many comparisons. Uniform cell decomposition
therefore gives finitely many monotone arcs of the one-step backward
singularity set, uniformly in the parameters. This decomposition is
made in the trigonometric charts, before normalized arclength is used.

On a fixed candidate region the entry time is uniformly one-half
H\"older up to its tangency boundary. Away from tangency this is the
implicit function theorem with a nonzero transverse derivative.
At tangency, a defining function for the boundary has a nonzero
second derivative along the tangent flight direction, bounded below
by the positive curvature bound. The parameter-dependent quadratic
normal form at this nondegenerate contact writes the two nearby
intersection times as a smooth center plus or minus a smooth bounded
factor times the square root of a nonnegative parameter function.
A finite cover of the compact tangency set makes these bounds uniform.
The square-root inequality gives the asserted exponent. The first
admissible root is the entry root; a zero root at the current obstacle
is discarded. The displacement label is constant on each candidate
region. Strict convexity precludes a higher-order tangency in this
one-flight calculation.

Backward singularity arcs lie in the unstable cone, uniformly
transverse to stable test curves. Their finitely many monotone pieces
therefore have uniformly bounded stable intersection counts and
$O(\epsilon)$ stable length in boundary neighborhoods. The smooth
monotone arclength change preserves these estimates. This verifies
the preceding-flight multiplier assumptions without a long-itinerary
partition. Regular matched connectors and their intermediate
singularity cuts are those of the unweighted collision construction,
so the action proof introduces no new matching partition.

For the fixed convex polygonal observation cell, a flight has finitely
many cell intervals. Exclude grazing and collision singularities,
cell-edge tangencies and vertex crossings. These are finitely many
semialgebraic curves and boundary pieces with zero collision measure;
line coincidence with an edge is included among the exceptional
configurations. On compact sets separated from them the intervals and
their endpoints vary uniformly with the parameters. Closed neighborhoods
of the limiting exceptional sets can be chosen with arbitrarily small
measure and null boundary, by taking continuity radii. These are the
closed exceptional neighborhoods required in the corrected varying-test
statement. Thus the observation hypothesis, not just pointwise
convergence away from an open set, is verified.

All covariance and normalization continuity statements now follow
from the compact-family principle. Its finite-cover arithmetic gives
uniform positivity by compactness. To compute the normalizations,
integrate $h+h''$ for the perimeter and use
$|O|=\tfrac12\int_0^{2\pi}(h^2-(h')^2)\,d\psi$ for the area.
Orthogonality of the third and fourth harmonics gives
\eqref{eq:v37-support-area}; the mean roof is the same invariant
collision/flow normalization as in the ellipse application.

Finally a centrally symmetric body's support function, after a
translation, is $\pi$-periodic; the translation contributes only a
first harmonic. The nonzero third harmonic in \eqref{eq:v37-support-family}
cannot be removed by such a translation. For $e\ne0$ the obstacle is
therefore not centrally symmetric. Parameter redundancy when an
amplitude vanishes does not singularize the geometry or the covariance.
\end{proof}

\paragraph{The class of selected conclusions.}
Regular endpoint selectors always mean the null-boundary and closed
exceptional-neighborhood class of the compact-family theorem. The
product of their stationary means must be uniformly positive for a
selected denominator. The arbitrary bounded trajectory-test conclusion
concerns only the comparison of posterior measures, not a Gaussian
amplitude for an arbitrary path selector. We use
$\|\mu-\nu\|_{\TV}=\sup_{|F|\le1}|\mu(F)-\nu(F)|$;
for probabilities the often used distance
$\sup_A|\mu(A)-\nu(A)|$ is one half of this norm. All comparison
bounds hold in this explicitly chosen, larger convention.
